% Alles was im Kopf steht, ausser Funktionsdefinitionen. Diese sollten besser in
% eine eigene Datei ausgelagert werden. Fuer kurze Funktionen einfach in
% functions.tex, sonst ggf. fuer zusammenhaengende Dinge eine eigene Datei.

% Titel, Autor etc. sollten gesetzt werden. Man kann sie bei Bedarf einfach in
% der jeweiligen Praesentation ueberschreiben.

\usepackage[ngerman]{babel}
\usepackage[utf8]{inputenc}
\usepackage[T1]{fontenc}
\usepackage{etoolbox}
\usepackage{csquotes}
\usepackage{amsmath,amssymb,amsthm}
\usepackage{color}
\usepackage{textcomp}
\usepackage{mathtools}
\usepackage[locale=DE]{siunitx}
\usepackage{xr}
\usepackage[version=3,arrows=pgf]{mhchem}
\usepackage{lipsum}
\usepackage{wrapfig}
\usepackage{fancyvrb}

\usepackage{tikz}
% Für lustige Grafiken
\usetikzlibrary{shapes,arrows}
\usetikzlibrary{positioning}
\usetikzlibrary{arrows}
\usetikzlibrary{calc,fadings,decorations.pathreplacing}

\usepackage[european]{circuitikz}

\hypersetup{
    bookmarks=true,
    unicode=false,
    pdfstartview={FitH},
    pdftitle={Eine Einführung in LaTeX},
    pdfauthor={Arbeitsgruppe LaTeX, Universität Bremen},
    hidelinks=true
}

\definecolor{coolgray}{rgb}{0.90625,0.8828125,0.8359375}
\definecolor{BurntOrange}{rgb}{1,0.5,0}
\definecolor{Brown}{HTML}{C17D11}
\definecolor{SeaGreen}{HTML}{2e8b57}

\graphicspath{{./images/}}

\renewcommand{\familydefault}{\sfdefault}
\usetheme[footlinenumber, navline=false, footlineauthor]{Bremen}

\fbnum{00}
\fbname{Uni Bremen}
\subject{\LaTeX}

\title{Eine Einführung in \LaTeX{}}
\author[AG \LaTeX]{Arbeitsgruppe \LaTeX\\\smallskip
\footnotesize{}Robin Görmer, Malte Gerken\\
Yannik Schädler}
\date{Orientierungswoche \the\year{}}


\institute[Uni Bremen]{-- Universität Bremen --}

\makeatletter
% Titelseite automatisch erzeugen
\g@addto@macro\beamer@lastminutepatches{\begin{frame}
    \titlepage{}
\end{frame}}

% Keine Abstände vor und nach Codeblöcken
\preto{\@verbatim}{\topsep=1pt \partopsep=1pt }
\makeatother

\AtBeginSection{\begin{frame}
    \begin{center}
    \Huge \textbf{\secname}
    \end{center}        
\end{frame}}

\makeatletter
\setbeamertemplate{footline}{%
  \usebeamerfont{subsection in head/foot}%
  \mbox{}\rlap{%
    \begin{pgfpicture}{0pt}{0pt}{\paperwidth}{2.8\baselineskip}%
      % Grauer Fond
      \color{coolgray 1}%
      \pgfpathrectangle{\pgfpoint{0cm}{0cm}}{%
	\pgfpoint{\paperwidth}{1.6\baselineskip}}\pgfusepath{fill}
      \pgfpathrectangle{\pgfpoint{0mm}{0cm}}{%
	\pgfpoint{.75\paperwidth}{2.8\baselineskip}}\pgfusepath{fill}
      \color{black}%
      % Universitaetslogo
%      \pgfputat{\pgfpoint{4mm}{1.5\baselineskip}}{%
%	\pgfbox[left,bottom]{\pgfimage[height=2\baselineskip]{unilogo}}%
      %}
    \end{pgfpicture}%
  }
  %% Universitätslogo
  \rlap{%
    \raisebox{.55\baselineskip}{%
      \mbox{}\hspace{4mm}%
      \parbox[b]{.75\paperwidth}{%
	\href{http://www.uni-bremen.de}{%
	  \includegraphics[height=1.7\baselineskip]{unilogo}}%
	}}}%
  %% Institute/Partnerlogos
  \rlap{%
    \raisebox{1.75\baselineskip}{%
      \mbox{}\hspace{.5\paperwidth}%
      \llap{%
      \parbox[t]{.15\paperwidth}{%
        \vspace{0pt}%
	\mbox{}\hfill%
	\beamer@theme@footline@partnerlogo}%
	\hspace{.025\paperwidth}}
      \rlap{%
      \hspace{.025\paperwidth}%
      \parbox[t]{0.15\paperwidth}{%
        \vspace{0pt}%
	\beamer@theme@footline@secondarypartnerlogo}%
	}}}%
  %% Autor
  \rlap{%
    \raisebox{1.9\baselineskip}{%
      \mbox{}\hspace{.75\paperwidth}%
      \parbox[b]{.25\paperwidth}{%
	\mbox{}\hspace{.5mm}\beamer@theme@footline@author\hfill%
% Ausgelager siehe unten
%	\beamer@theme@footline@number\hfill~\hspace{5mm}
  }}}%
  %%% Foliennummer
  \rlap{%
    \raisebox{1.9\baselineskip}{%
      \mbox{}\hspace{.91\paperwidth}%
      \parbox[b]{.05\paperwidth}{%
        \raggedright\beamer@theme@footline@number
      }
    }
  }
  %% Abschnittsleiste
  \rlap{\raisebox{1.9\baselineskip}{%
      \beamer@theme@footline@section%
  }}%
  %% Navigationssymbole
  \rlap{\parbox[b]{\paperwidth}{%
      \hfill\usebeamertemplate***{navigation symbols lower}%
      \hspace{1.5mm}\mbox{}\vspace{.5mm}%
  }}%
}
\makeatother



\newenvironment{fframe}{\begin{frame}[fragile,environment=specialframe]}{\end{frame}}


\makeatletter
        \providecommand\grpn[1]{\textcolor{OliveGreen}{\texttt{#1}}}
        \providecommand\cls[1]{\textcolor{BurntOrange}{\textsf{#1}}}
        \providecommand\pkg[1]{\cls{#1}}
        \providecommand\umg[1]{\textcolor{SeaGreen!50!Brown}{\texttt{#1}}}
        \providecommand\msur[1]{\texttt{\{#1\}}}
        \providecommand\osur[1]{\texttt{[#1]}}
        \providecommand\meta[1]{\emph{\ensuremath\langle#1\ensuremath\rangle}}
% \cmd{\foo} Prints \foo verbatim
        \def\cmd#1{\cs{\expandafter\cmd@to@cs\string#1}}
        \def\cmd@to@cs#1#2{\char\number`#2\relax}
        \DeclareRobustCommand\cs[1]{\texttt{\char`\\#1}}
% \marg \marg{text} prints {text}, `mandatory argument'.
        \providecommand\marg[1]{%
        {\ttfamily\char`\{}\meta{#1}{\ttfamily\char`\}}}
% \oarg \oarg{text} prints [text], `optional argument'.
        \providecommand\oarg[1]{%
        {\ttfamily[}\meta{#1}{\ttfamily]}}
% \parg \parg{te,xt} prints (te,xt), `picture mode argument'.
        \providecommand\parg[1]{%
        {\ttfamily(}\meta{#1}{\ttfamily)}}

% Keine Abstände vor und nach Codeblöcken
\preto{\@verbatim}{\topsep=1pt \partopsep=1pt }
\makeatother

\newcommand{\uebung}[1]{
\begin{frame}
\frametitle{Zeit zum Ausprobieren}

Die Folien findet ihr hier: \url{\slideurl}

Die Aufgaben hier: \url{\exerciseurl}
\bigskip
\begin{center}
    {\huge
    Bearbeitet jetzt die\\
    Aufgabe #1
    }
\end{center}
\end{frame}}
